%! The Complete Solution to Ax = b — Unit 3, Lesson 3.3 — Slides (Beamer)
\documentclass[aspectratio=169, 11pt]{beamer}
\usepackage{linalg-beamer}   % provides \burgundyheader{} and \sectionlabel[color]{}
\newcommand{\vv}{\mathbf{v}}
\newcommand{\ww}{\mathbf{w}}
\newcommand{\xx}{\mathbf{x}}
\newcommand{\yy}{\mathbf{y}}
\newcommand{\bb}{\mathbf{b}}
\newcommand{\svec}{\mathbf{s}}
\newcommand{\zero}{\mathbf{0}}

\begin{document}

% TITLE SLIDE (CourseName is not defined in beamer, so the name is literal)
{
\setbeamercolor{background canvas}{bg=burgundy}
\begin{frame}[plain]
  \vspace{2.8cm}\hspace{0.55cm}%
  \begin{minipage}{0.9\textwidth}
    {\color{goldacc}\bfseries\small LESSON 3.3}\\[0.5cm]
    {\color{white}\bfseries\LARGE The Complete Solution to $A\xx=\bb$}\\[1.0cm]
    {\color{blushmid}\small Linear Algebra~$\cdot$~Unit 3: The Four Fundamental Subspaces}
  \end{minipage}
\end{frame}
}

% HOOK
\begin{frame}[t, plain]
  \burgundyheader{Which settings hit the target $\bb$?}
  \sectionlabel{Hook}
  \begin{columns}[T]
    \begin{column}{0.58\textwidth}
      The robot arm's tip moves to $A\xx$. In 3.2 we found the \emph{do-nothing} settings ($A\xx=\zero$).
      Now aim at a target $\bb$.\\[8pt]
      \textbf{Find \emph{one} winning setting} $\xx_p$ ($A\xx_p=\bb$).\\[4pt]
      Then add \emph{any} do-nothing move $\xx_n$:
      \[ A(\xx_p+\xx_n)=\bb+\zero=\bb. \]
      Still on target! So the winners are $\xx_p$ + the whole nullspace.
    \end{column}
    \begin{column}{0.40\textwidth}
      \centering
      \begin{tikzpicture}[scale=0.5]
        \draw[step=1, linegray!45, thin] (-2.3,-1.3) grid (2.3,3.3);
        \draw[->, thick, charcoal] (-2.3,0)--(2.5,0);
        \draw[->, thick, charcoal] (0,-1.3)--(0,3.5);
        \draw[very thick, burgundy] (-2,-1)--(2,1) node[right]{\scriptsize $N(A)$};
        \draw[very thick, royalblue, dashed] (-2,1)--(2,3);
        \draw[->, very thick, goldacc] (0,0)--(0,2);
        \fill[royalblue] (0,2) circle (2.5pt) node[right]{\scriptsize $\xx_p$};
        \fill[black] (0,0) circle (2.5pt);
      \end{tikzpicture}
    \end{column}
  \end{columns}
  \vspace{2pt}
  \begin{block}{The complete solution}
    Every solution of $A\xx=\bb$ is one particular solution plus the whole nullspace.
  \end{block}
\end{frame}

% BIG IDEA: PARTICULAR + NULLSPACE
\begin{frame}[t, plain]
  \burgundyheader{Particular $+$ nullspace $=$ complete}
  \sectionlabel{The big idea}
  If $A\xx_p=\bb$ and $A\xx_n=\zero$, then:
  \begin{itemize}
    \setlength{\itemsep}{5pt}
    \item \textbf{Still a solution:} $A(\xx_p+\xx_n)=A\xx_p+A\xx_n=\bb+\zero=\bb$.
    \item \textbf{Nothing else:} if $A\xx=\bb$ too, then $A(\xx-\xx_p)=\zero$, so $\xx-\xx_p\in N(A)$.
  \end{itemize}
  \vspace{4pt}
  \begin{block}{Complete solution}
    \[ \xx=\xx_p+\xx_n \qquad(\text{one particular solution} \;+\; \text{the whole nullspace}). \]
  \end{block}
\end{frame}

% FIND x_p FROM AUGMENTED MATRIX
\begin{frame}[t, plain]
  \burgundyheader{Find $\xx_p$: reduce $[A\mid\bb]$}
  \sectionlabel[goldacc]{Particular solution}
  Attach $\bb$ as an extra column and reduce; then set the free variable to $0$:
  \[
    \left[\begin{array}{ccc|c} 1 & 1 & 2 & 3 \\ 1 & 2 & 3 & 5 \end{array}\right]
    \longrightarrow
    \left[\begin{array}{ccc|c} 1 & 0 & 1 & 1 \\ 0 & 1 & 1 & 2 \end{array}\right].
  \]
  \begin{itemize}
    \setlength{\itemsep}{4pt}
    \item Columns $1,2$ pivot; column $3$ is \textbf{free} --- set $x_3=0$.
    \item Read back: $x_1=1,\ x_2=2$, so $\xx_p=\begin{bsmallmatrix} 1\\2\\0 \end{bsmallmatrix}$.
    \item Check: $1\begin{bsmallmatrix} 1\\1 \end{bsmallmatrix}+2\begin{bsmallmatrix} 1\\2 \end{bsmallmatrix}=\begin{bsmallmatrix} 3\\5 \end{bsmallmatrix}=\bb.\ \checkmark$
  \end{itemize}
\end{frame}

% COMPLETE SOLUTION + PICTURE
\begin{frame}[t, plain]
  \burgundyheader{The complete solution --- shifted nullspace}
  \sectionlabel{Anchor $+$ wiggle room}
  \begin{columns}[T]
    \begin{column}{0.56\textwidth}
      Same $A$ as 3.2: nullspace is all $t\,\svec$ with $\svec=\begin{bsmallmatrix} -1\\-1\\1 \end{bsmallmatrix}$. So
      \[ \xx=\begin{bmatrix} 1\\2\\0 \end{bmatrix}+t\begin{bmatrix} -1\\-1\\1 \end{bmatrix}. \]
      \begin{itemize}
        \setlength{\itemsep}{3pt}
        \item Only the nullspace part carries $t$; $\xx_p$ is a fixed anchor.
        \item The solution set is $N(A)$ \textbf{slid} through $\xx_p$ --- parallel, off the origin.
        \item Through the origin only when $\bb=\zero$.
      \end{itemize}
    \end{column}
    \begin{column}{0.42\textwidth}
      \centering
      \begin{tikzpicture}[scale=0.5]
        \draw[step=1, linegray!45, thin] (-2.3,-1.8) grid (2.3,3.3);
        \draw[->, thick, charcoal] (-2.3,0)--(2.5,0);
        \draw[->, thick, charcoal] (0,-1.8)--(0,3.5);
        \draw[very thick, burgundy] (-2,-1)--(2,1) node[right]{\scriptsize $N(A)$};
        \draw[very thick, royalblue, dashed] (-2,1)--(2,3) node[right]{\scriptsize $\xx_p{+}N(A)$};
        \draw[->, very thick, goldacc] (0,0)--(0,2);
        \fill[royalblue] (0,2) circle (2.5pt) node[left]{\scriptsize $\xx_p$};
        \fill[black] (0,0) circle (2.5pt);
      \end{tikzpicture}
    \end{column}
  \end{columns}
\end{frame}

% SOLVABILITY
\begin{frame}[t, plain]
  \burgundyheader{Is $A\xx=\bb$ even solvable?}
  \sectionlabel[goldacc]{Consistency}
  A solution exists exactly when $\bb$ lies in $C(A)$ (a reachable combination of the columns).
  \[
    \left[\begin{array}{cc|c} 1 & 1 & 1 \\ 2 & 2 & 3 \end{array}\right]
    \xrightarrow{\;R_2-2R_1\;}
    \left[\begin{array}{cc|c} 1 & 1 & 1 \\ 0 & 0 & 1 \end{array}\right]
    \;\Rightarrow\; 0=1 \text{ (no solution).}
  \]
  \begin{itemize}
    \setlength{\itemsep}{4pt}
    \item \textbf{Solvability:} every zero row of $A$ must meet a zero on the right.
    \item Full column rank $r=n$: $0$ or $1$ solution. \quad Full row rank $r=m$: always solvable.
    \item Invertible square ($r=m=n$): \emph{exactly one} for every $\bb$; a free variable $\Rightarrow$ infinitely many.
  \end{itemize}
\end{frame}

% ROADMAP
\begin{frame}[t, plain]
  \burgundyheader{Where this goes}
  \sectionlabel{Connections}
  \textbf{Today:} the complete solution $\xx=\xx_p+\xx_n$ --- reduce $[A\mid\bb]$ for $\xx_p$, add the
  whole nullspace, read the shifted geometry, and test solvability.\\[8pt]
  \begin{itemize}
    \setlength{\itemsep}{4pt}
    \item \textbf{3.4} Independence, Basis, and Dimension: which special solutions are really independent.
    \item \textbf{3.5} Dimensions of the Four Subspaces.
    \item \textbf{3.6} The Fundamental Theorem of Linear Algebra.
  \end{itemize}
  \vspace{4pt}
  \begin{block}{Keep in mind}
    The solution set is a subspace \emph{only} when $\bb=\zero$; otherwise it is $N(A)$ shifted off the origin.
  \end{block}
\end{frame}

\end{document}
